\begin{tikzpicture}[
        base/.style={rectangle, draw, minimum width=2.5cm, minimum height=1cm, text=white, rounded corners, text width=2.8cm, align=center, inner sep=8pt},
        server/.style={base, fill=RoyalBlue},
        client/.style={base, fill=ForestGreen},
        arrow/.style={-Stealth, thick, rounded corners},
        code/.style={font=\ttfamily\small},
        socket/.style={circle, draw, fill=yellow!10, minimum size=0.5cm}
    ]

    % Socket symbol inside server rectangle
    \node[client] (client) at (0, 0) {Client\\\texttt{198.168.0.10}};
    \node[server] (server) at (10, 0) {Server\\\texttt{198.168.0.11}};
    % State boxes
    \node[rectangle, draw, fill=gray!10, below=3cm of client.east, text width=7cm, align=left, font=\small] (client_state) {
        \textbf{Client:} \\
        \iftoggle{beamer}{
            \only<4->{1. Socket erstellen (\texttt{.socket()})\\}
            \only<5->{2. Mit Server-Socket verbinden (\texttt{.connect()})\\}
            \only<6->{3. Nachricht senden (\texttt{.send()})\\}
            \only<9->{4. Antwort empfangen (\texttt{.recv()})}
        }{
            1. Socket erstellen (\texttt{.socket()})\\
            2. Mit Server-Socket verbinden (\texttt{.connect()}) \\
            3. Nachricht senden (\texttt{.send()})\\
            4. Antwort empfangen (\texttt{.recv()})
        }
    };

    \node[rectangle, draw, fill=gray!10, below=3cm of server.east, xshift=-2cm, text width=7cm, align=left, font=\small] (server_state) {
        \textbf{Server:} \\
        \iftoggle{beamer}{
            \only<1->{1. Socket erstellen (\texttt{.socket()})\\}
            \only<2->{2. Socket binden (\texttt{.bind()})\\}
            \only<3->{3. Warten auf Verbindungen (\texttt{.listen()})\\}
            \only<5->{4. Accept Verbindung (\texttt{.accept()})\\}
            \only<7->{5. Nachricht empfangen (\texttt{.recv()}) \\}
            \only<8->{6. Antwort senden (\texttt{.send()})}
        }{
            1. Socket erstellen (\texttt{.socket()})\\
            2. Socket binden (\texttt{.bind()})\\
            3. Warten auf Verbindungen (\texttt{.listen()})\\
            4. Accept Verbindung (\texttt{.accept()})\\
            5. Nachricht empfangen (\texttt{.recv()}) \\
            6. Antwort senden (\texttt{.send()})
        }

    };
    \node[socket] (server_socket_placeholder) at (server.north west) {};
    \iftoggle{beamer}{\pause}{}

    \node[socket] (server_socket) at (server.north west) {\large \emoji{pig-nose}};

    % Arrows with labels
    \node[rectangle, draw, fill=yellow!5, above right=.2cm and .6cm of server_socket, minimum width=2.2cm, minimum height=1.2cm, font=\small, align=center, rounded corners] (server_socket_text){
        \textbf{Socket:}\\
        IP: \texttt{198.168.0.11}\\
        Port: \texttt{12345}\\
        Protokoll: \texttt{TCP}
    };

    % draw dotted line from (server_socket) to (server_socket_text)
    \draw[dashed] (server_socket) -- (server_socket_text);

    \iftoggle{beamer}{\pause}{}
    \draw[arrow, loop, min distance=2cm, in=120, out=60] (server_socket.north) to node[above, code, align=center, text width=2cm] {bind()\\ listen()\\ \iftoggle{beamer}{\only<5->{accept()}}{accept()}} (server_socket.north);

    \iftoggle{beamer}{\pause}{}
    \node[socket] (client_socket) at (client.north east) {};
    \iftoggle{beamer}{\pause}{}
    \draw[arrow] (client_socket) to[bend left=40] node[above, yshift=3pt, code] {connect() \faPlug{}} (server_socket);
    \iftoggle{beamer}{\pause}{}
    \draw[arrow, bend left=10] (client_socket) to node[above, code] {send("Hallo Server...")} (server_socket);
    \iftoggle{beamer}{\pause}{}
    \draw[arrow, dashed, bend left=10] (server_socket) to node[below, yshift=3pt, code, xshift=10pt] {recv()} (client_socket);
    \iftoggle{beamer}{\pause}{}
    \draw[arrow, dashed, bend left=40] (server_socket.south) to node[below, yshift=-8pt, code, xshift=-10pt] {send("Hallo Client...")} (client_socket.south);

\end{tikzpicture}